\chapter{Introduction}

\section{Background and motivation}
Restaurant reviews contain detailed descriptions of food, service, price, atmosphere and
accessibility. Conventional restaurant recommenders commonly predict ratings or popularity,
but a user often expresses a need in natural language: ``a quiet place with outdoor seating'',
``affordable food'', or ``something sour and spicy''. Review-based retrieval can answer this
need directly by finding review evidence that matches the query and using it to rank restaurants.

The central motivation of this work is explainability. A ranked restaurant without evidence is
difficult for a user to verify. Our system returns existing review sentences alongside each
restaurant, allowing a user to inspect why an item was retrieved. This is different from
generating an explanation after recommendation: the explanation is selected from the same
review evidence used by the retrieval system.

The project began with a broader Thai food recommendation proposal. The data audit showed that
the English RIRD benchmark contains restaurant identities and relevance judgements, while the
available Wongnai data contains Thai reviews but no restaurant identities. We consequently
separated the claims: English supports measured restaurant ranking and explainability; Thai
supports a sentence-level transfer case study. This change prevents a common evaluation error
in which a model is described as a Thai restaurant recommender without a Thai restaurant-level
ground truth.

\section{Business understanding and potential impact}
The intended users are diners choosing a restaurant from a personal preference rather than a
fixed category. The practical benefits are:
\begin{itemize}
  \item natural-language search instead of manually browsing categories;
  \item transparent recommendations supported by real review text;
  \item better handling of specific constraints such as price, waiting time and atmosphere;
  \item multilingual sentence retrieval as a foundation for Thai food discovery.
\end{itemize}
For restaurant platforms, this can increase discovery of suitable businesses and reduce the
trust gap caused by opaque ranking. The system does not claim to replace human judgement:
reviews can be subjective, outdated or contradictory. Its value is to make relevant evidence
easier to find.

The business problem is not simply ``predict the highest-rated restaurant''. A user may accept
a lower-rated restaurant because it is inexpensive, quiet, accessible or suitable for a dietary
constraint. A global rating is therefore an incomplete proxy for utility. Query-conditioned
review retrieval gives the platform a way to serve long-tail preferences while giving the user
enough evidence to decide whether the recommendation is credible. For restaurants, this may
improve discovery when a useful characteristic is present in reviews but not in a category
label.

\section{Stakeholders and success criteria}
The primary stakeholder is the diner, who needs relevant choices and understandable evidence.
Secondary stakeholders are restaurant platforms, which need a retrieval service that is
measurable, computationally practical and robust to uneven review volume. We define success as:
\begin{enumerate}
  \item better query-conditioned ranking than a strong lexical baseline;
  \item evidence that is traceable to original reviews;
  \item acceptable behaviour for vague and negative queries;
  \item a complexity level suitable for a small interactive demo.
\end{enumerate}

\section{Research objectives}
The project has three objectives:
\begin{enumerate}
  \item compare BM25 and multilingual dense retrieval fairly using the same review units and
  aggregation functions;
  \item evaluate whether retrieved sentences are faithful to ranking and supportive of a query;
  \item test sentence-level transfer to real Thai restaurant reviews using independently
  annotated taste queries.
\end{enumerate}

\section{Contributions}
The contributions are an evidence-based reviewed-item retrieval pipeline, a controlled
comparison of review- and sentence-level retrieval, an analysis of review-volume bias, a
sentence-level explanation interface, and a reproducible Thai taste case study. The live demo
also includes query rewriting and selective cross-encoder reranking for difficult negation and
exclusion queries.
